\par\medskip
\noindent\textbf{Corolario 5.7.}
Sea $S$ un esquema formal noetheriano y $\mathfrak X$ un esquema
formal $S$-ádico esencialmente de tipo finito. Entonces:

\smallskip
\noindent\textbf{(1)} La formación de
$J_{\mathfrak X}^p(\mathfrak X|S)$ conmuta con el cambio de base. Más
precisamente, si
\[
\mathfrak X'=\mathfrak X\times_SS',
\]
entonces
\[
J_{\mathfrak X'}^p(\mathfrak X'|S')
=J_{\mathfrak X}^p(\mathfrak X|S)\times_SS'.
\]

\smallskip
\noindent\textbf{(2)} Las componentes conexas de
$J_{\mathfrak X}^p(\mathfrak X|S)$ están en correspondencia biyectiva
con las de $J_{X_0}^p(X_0|S_0)$. Para todo
\[
x\in\bigl|J_{X_0}^p(X_0|S_0)\bigr|
\]
hay un isomorfismo canónico
\[
J_{\mathfrak X}^p(\mathfrak X|S)_{(x)}
\simeq
J_{\mathfrak X_{(x)}}^p(\mathfrak X_{(x)}|S).
\]
En particular, si
\[
\bigl|J_{X_0}^p(X_0|S_0)\bigr|=\{x_1,\ldots,x_m\}
\]
es discreto, entonces
\[
J_{\mathfrak X}^p(\mathfrak X|S)
=\coprod_{1\leq i\leq m}
 J_{\mathfrak X_{(x_i)}}^p(\mathfrak X_{(x_i)}|S)
\qquad\text{(unión disjunta)}.
\]

%%%%%%%%%%  scan idx 132  |  folio 121  %%%%%%%%%%

\smallskip
\noindent\textbf{(3)} Sea $S$ un esquema noetheriano ordinario y
$X$ un $S$-esquema esencialmente de tipo finito. Dado un subesquema
cerrado $S_0\subset S$, sea $\widehat S$ la completación formal de
$S$ a lo largo de $S_0$. Si
\[
\widehat X=X\times_S\widehat S
\]
es la completación formal de $X$ a lo largo de
$X_0=X\times_SS_0$, entonces
\[
J_{\widehat X}^p(\widehat X|\widehat S)
=J_X^p(X|S)\times_S\widehat S.
\]
Equivalentemente, es la completación formal de $J_X^p(X|S)$ a lo
largo de $J_{X_0}^p(X_0|S_0)$.

\smallskip
\noindent\textbf{(4)} Supongamos que un entero fijo $n\geq0$
satisface la condición 5.4(1) o 5.4(2) para todo
$X_\nu\longrightarrow S_\nu$. Entonces
\[
x\notin J_{\mathfrak X}^n(\mathfrak X|S)
\]
si y solo si $\mathfrak X$ es formalmente liso sobre $S$ en $x$.

\smallskip
\noindent\textit{Demostración.}
La afirmación (1) se deduce de
\[
\Omega_{\mathfrak X'|S'}
\simeq
\mathcal O_{S'}\widehat\otimes_{\mathcal O_S}
\Omega_{\mathfrak X|S},
\]
mientras que (2) se deduce de
\[
\mathcal O_{\mathfrak X_{(x)}}
\widehat\otimes_{\mathcal O_{\mathfrak X}}
\Omega_{\mathfrak X|S}
\simeq\Omega_{\mathfrak X_{(x)}|S}.
\]
Para (3),
\[
\begin{aligned}
J_{\widehat X}^p(\widehat X|\widehat S)
&=\varinjlim_\nu
 J_{X\times_SS_\nu}^p(X\times_SS_\nu|S_\nu) \\
&=\varinjlim_\nu\bigl(J_X^p(X|S)\times_SS_\nu\bigr)
 =J_X^p(X|S)\times_S\widehat S.
\end{aligned}
\]
Finalmente, como $\mathfrak X$ es $S$-ádico,
\[
\begin{aligned}
\mathfrak X\to S\text{ es formalmente liso en }x
&\Longleftrightarrow
 X_\nu\to S_\nu\text{ es formalmente liso en }x
 \text{ para todo }\nu \\
&\Longleftrightarrow
 x\notin J_{X_\nu}^n(X_\nu|S_\nu)
 \text{ para todo }\nu \\
&\Longleftrightarrow
 x\notin\varinjlim_\nu J_{X_\nu}^n(X_\nu|S_\nu)
 =J_{\mathfrak X}^n(\mathfrak X|S).
\end{aligned}
\]
\hfill$\square$
